\section*{Arithmetic resolution and quantum refinement}
The preceding constructions compare geometric and spectral resolutions. Arithmetic refinement provides another concrete family: finite unit registers, Haar packets, hierarchical Hamiltonians, and polyphase channels. The question remains the same. Which maps preserve the state, which preserve the dynamics, and what additional structure appears when hidden components are eliminated?

Normalized Haar averages identify coarse states inside a fine Hilbert space. Compatibility of the Hamiltonian, rather than normalization alone, determines whether that subspace carries the same evolution. Schur reduction and its time-domain memory equation describe the effect of components that cannot simply be discarded.

Infinite depth introduces a separate analytic problem. The hierarchical model specifies a common Hilbert space, a maximal operator domain, and the topology of approximation. Spectral decimation and hypergeometric companion functions provide exact coordinates within their declared domains. None of these identities alone identifies arithmetic depth with a physical observable.

The following account of covariant fiber dynamics supplies a candidate physical coupling and examines its energy and production constraints. That transition separates a mathematical refinement structure from a proposed model of matter.

The visual and verification appendices distinguish exact identities, finite computations, and formal coverage. The infinite arithmetic hierarchy is presented with its written proofs and is not assigned a Lean certificate beyond the statements specifically identified there.

\clearpage
\section*{Reading routes and conventions}
\addcontentsline{toc}{section}{Reading routes and conventions}
The finite route begins with unit coordinates and Haar normalization. It continues through exact intertwining and controlled leakage. The spectral route begins with the block resolvent and follows radial chains, residue channels, and Chebyshev elimination. The limiting route begins with cylinder projections, constructs the maximal hierarchical Hamiltonian, and then treats compatible tensor dynamics.

An embedding and an elimination have different jobs. An isometry retains the old inner product. A unitary polyphase transformation retains the full fine operator in a new basis. A Schur complement solves for hidden coordinates on a regular spectral chart. A projection approximation converges only in the topology established by its proof. These operations should be identified before their consequences are compared.

\begingroup\small\renewcommand{\arraystretch}{1.1}
\begin{longtable}{@{}L{.24\linewidth}L{.67\linewidth}@{}}
\toprule
Convention & Meaning\\\midrule\endhead
Hilbert normalization & Finite registers use an explicitly stated orthonormal basis. Cylinder functions use probability Haar measure. The conversion between counting coordinates and normalized indicators includes the square root of the cylinder measure.\\
Refinement & An isometry $V$ has $V^*V=I$. Its old-range projector is $P=VV^*$ and its new-detail projector is $Q=I-P$. Exact dynamical compatibility is a separate identity, $H'V=VH$.\\
Hierarchy & The finite chapter introduces cumulative projectors $E_r$ and their details. The infinite chapter writes $\mathsf E_m$ for conditional expectations and $\mathsf F_m=\mathsf E_m-\mathsf E_{m-1}$ for depth details, with the constant sector declared separately.\\
Self-adjointness & A finite Hamiltonian is Hermitian. Complex coupling amplitudes appear with conjugates in the reverse direction. An unbounded Hamiltonian also requires a dense domain and a proof that the declared realization is self-adjoint.\\
Spectral parameter & Schur formulas use $z-H$ and require the hidden resolvent $(z-QHQ)^{-1}$. Real-energy formulas are restricted to their regular chart. They are not evaluated at hidden poles without an extension argument.\\
Spectral coordinates & A residue or polyphase transformation is a change of basis. Selecting a channel, tracing a subsystem, averaging amplitudes, and eliminating hidden coordinates are different operations.\\
Limits & Strong vector, strong operator, graph-norm, and operator-norm convergence have distinct quantifiers. The infinite chapter proves the stated limits for its compatible block model.\\
Evidence & Ordinary derivations, fresh bounded calculations, retained program outputs, and archived formal source are recorded separately. A small residual does not prove a universal infinite-depth or physical claim.\\
\bottomrule
\end{longtable}
\endgroup

The technical source concordance identifies statements by exact source member, formula, and claim. Each result here carries its own hypotheses. A source title or a recorded test result is not a substitute for those hypotheses.

\clearpage
